\documentclass[conference,compsoc]{IEEEtran}
\usepackage{amsmath,amsfonts,amssymb}
\usepackage{array}
\usepackage[caption=false,font=footnotesize,labelfont=sf,textfont=sf]{subfig}
\usepackage{textcomp}
\usepackage{stfloats}
\usepackage{url}
\usepackage{verbatim}
\usepackage{graphicx}
\usepackage[nocompress]{cite}
\usepackage{listings}
\usepackage{xcolor}
\usepackage[catalan]{babel}
\usepackage[T1]{fontenc}
\usepackage[utf8]{inputenc}
\usepackage{booktabs}
\usepackage{float}
\usepackage{cuted}
\usepackage{wrapfig}
\usepackage{multirow}
\usepackage{microtype}
\usepackage{siunitx}
\usepackage{pgfplots}
\pgfplotsset{compat=1.18}
\graphicspath{{./}}

\renewcommand{\abstractname}{Resum}
\renewcommand{\IEEEkeywordsname}{Termes d'índex}

\makeatletter
\def\abstract{\normalfont\@IEEEtweakunitybaselinestretch{1.15}%
    \if@twocolumn
      \@IEEEabskeysecsize\noindent\textbf{\textit{\abstractname}}---\normalfont\@IEEEabskeysecsize\relax
    \else
      \bgroup\par\addvspace{0.5\baselineskip}\centering\vspace{-1.78ex}\@IEEEabskeysecsize\textbf{\abstractname}\par\addvspace{0.5\baselineskip}\egroup\quotation\normalfont\@IEEEabskeysecsize%
    \fi\@IEEEgobbleleadPARNLSP}
\def\IEEEkeywords{\normalfont\@IEEEtweakunitybaselinestretch{1.15}%
    \if@twocolumn
      \@IEEEabskeysecsize\vskip 0.5\baselineskip plus 0.25\baselineskip minus 0.25\baselineskip\noindent
      \textbf{\textit{\IEEEkeywordsname}}---\normalfont\@IEEEabskeysecsize\relax
    \else
      \bgroup\par\addvspace{0.5\baselineskip}\centering\@IEEEabskeysecsize\textbf{\IEEEkeywordsname}\par\addvspace{0.5\baselineskip}\egroup\quotation\normalfont\@IEEEabskeysecsize%
    \fi\@IEEEgobbleleadPARNLSP}
\makeatother

\lstset{
    language=Java,
    basicstyle=\ttfamily\footnotesize,
    keywordstyle=\color{blue}\bfseries,
    commentstyle=\color{gray}\itshape,
    stringstyle=\color{red!70!black},
    numbers=left,
    numberstyle=\tiny\color{gray},
    numbersep=5pt,
    frame=single,
    breaklines=true,
    showstringspaces=false,
    tabsize=4,
    captionpos=b
}

\hyphenation{Huff-man com-pres-sio des-com-pres-sio fre-quen-ci-es}

\begin{document}

\title{Compressió Huffman amb Modes Adaptatius\\ Pràctica 4}

\author{\IEEEauthorblockN{Serafí Nebot Ginard}
\IEEEauthorblockA{Universitat de les Illes Balears (UIB)\\
Grau en Enginyeria Informàtica\\
Algorismes Avançats, Curs 2025--2026}}

\maketitle

\begin{strip}
\begin{abstract}
Aquesta memòria descriu el disseny i la implementació d'una aplicació Java
per comprimir i descomprimir fitxers mitjançant codificació de Huffman.
El programa analitza la distribució de freqüències del fitxer, construeix
arbres de Huffman amb diferents estructures de cua de prioritat i tria
automàticament entre modes d'1 byte, 2 bytes i compressió per blocs. La
interfície gràfica permet seleccionar fitxers, visualitzar l'arbre generat,
consultar la taula de codis i executar comparatives de rendiment. També
s'ha implementat una CLI que reutilitza el mateix model que la GUI, de
manera que els benchmarks són reproduïbles sense interacció manual. Es
presenten proves unitàries centrades en comprimir i descomprimir sense pèrdua,
robustesa del format,
correcció de cues i exportació CSV, i una anàlisi
experimental sobre els fitxers de Silesia.
\end{abstract}

\begin{IEEEkeywords}
Huffman, compressió sense pèrdues, entropia, cua de prioritat, MVC,
Java Swing, benchmark, Silesia
\end{IEEEkeywords}
\end{strip}

\IEEEpeerreviewmaketitle

%% ================================================================
\section{Introducció}

La codificació de Huffman és una tècnica de compressió sense pèrdues:
el fitxer descomprimit ha de ser exactament igual que l'original. La idea
principal és senzilla: els símbols que apareixen més sovint reben codis
més curts, i els símbols menys freqüents reben codis més llargs. Aquests
codis formen un codi prefix, és a dir, cap codi complet és prefix d'un
altre; per això la descompressió pot llegir el flux de bits sense
separadors explícits entre símbols \cite{ref_gfg_huffman}.

L'objectiu de la pràctica no és només implementar Huffman, sinó construir
una aplicació completa al seu voltant. El programa ha de poder comprimir
fitxers de text o binaris, guardar un arxiu propi que es pugui
descomprimir en una execució futura, mostrar l'arbre i la taula de codis,
donar estadístiques de compressió i mesurar rendiment. A més, s'han
comparat diferents estructures de cua de prioritat, perquè la construcció
de l'arbre depèn directament d'extreure repetidament els dos nodes de
menor freqüència.

Una decisió important d'aquesta implementació és que la compressió no
produeix cap fitxer si l'arxiu resultant seria més gran que l'original.
Això fa que alguns fitxers petits o molt aleatoris fallin de forma
controlada: és preferible informar l'usuari que generar un resultat que
va en contra del propòsit del propi programa.

%% ================================================================
\section{Fonaments}

\subsection{Freqüències, Entropia i Codis}

Abans de comprimir, el programa compta quantes vegades apareix cada
símbol. Un símbol pot ser un byte individual o una parella de bytes,
segons el mode que s'estigui analitzant. Aquesta taula de freqüències és
la base de tot el procés: defineix quins símbols mereixen codis curts i
quins poden tenir codis més llargs.

L'entropia de Shannon dóna una estimació del nombre mínim de bits per
símbol que es podria esperar per a una distribució concreta:
\[
H = -\sum_i p_i \log_2(p_i)
\]
on $p_i$ és la probabilitat del símbol $i$. En aquesta pràctica s'utilitza
com a valor de referència i com a estimació ràpida en la selecció de mida
de bloc. No es modifica l'arbre directament amb l'entropia; l'arbre es
construeix amb l'algorisme de Huffman clàssic. La diferència és que el
programa calcula diversos plans possibles i selecciona el que promet una
mida final més petita.

La longitud mitjana del codi es calcula com:
\[
L = \frac{\text{bits totals codificats}}{\text{nombre de símbols}}
\]
Per als modes de 2 bytes, el símbol és una parella de bytes. Per això
els valors d'entropia i longitud mitjana s'han d'interpretar sempre dins
del mateix tipus de símbol; no és equivalent comparar directament
``bits per byte'' i ``bits per parella''.

\subsection{Construcció de l'Arbre de Huffman}

L'arbre es construeix de baix cap amunt. Primer es crea una fulla per a
cada símbol amb freqüència major a 0. Després es repeteix el mateix pas
fins que només queda un node: extreure els dos nodes de menor freqüència,
unir-los en un node intern i tornar-lo a inserir a la cua de prioritat.
El node final és l'arrel de l'arbre.

Un detall pràctic important és el desempat. Quan dos símbols tenen la
mateixa freqüència, hi pot haver diversos arbres òptims. La implementació
ordena també pel símbol mínim de cada subarbre. Això no millora la mida
de compressió, però fa que totes les estructures de cua generin el mateix
arxiu, cosa essencial per provar que les cues són equivalents.

\subsection{Diversos Modes de Compressió}

El fitxer es pot mirar amb diferents ``ulleres''. En mode d'1 byte, cada
valor entre 0 i 255 és un símbol. En mode de 2 bytes, cada parella
consecutiva és un símbol entre 0 i 65535. Si el fitxer conté patrons
repetits de dos bytes, aquest segon mode pot ser clarament millor.

També hi ha un mode per blocs. Un sol arbre global pot ser dolent quan
el fitxer canvia de comportament al llarg del temps. Per exemple, una
part pot estar dominada per un byte i una altra per un altre. El mode per
blocs divideix el fitxer i pren una decisió local per a cada bloc:
Huffman d'1 byte, Huffman de 2 bytes o bloc emmagatzemat sense comprimir.
El mode \texttt{STORED} només s'utilitza en aquest últim cas, com a
fallback local. No s'ofereix com a compressió de fitxer sencer perquè
només afegiria capçalera i faria créixer l'arxiu.

%% ================================================================
\section{Descripció de l'Aplicació}

L'aplicació té dues entrades: una interfície gràfica Swing i una línia
de comandes. Les dues reutilitzen les mateixes classes del model, de
manera que comprimir des de la GUI o des de la CLI executa el mateix
procés intern.

\subsection{Interfície Gràfica}

La finestra principal està organitzada en quatre pestanyes.

\begin{itemize}
    \item \textbf{Fitxers:} explorador integrat, ruta d'entrada, ruta de
          sortida, selector de cua de prioritat, selector de mode i resum
          d'estadístiques.
    \item \textbf{Taula de codis:} llista els símbols, valor hexadecimal,
          freqüència, probabilitat i codi assignat. Les columnes són
          ordenables, i les freqüències/probabilitats s'ordenen
          numèricament.
    \item \textbf{Arbre:} visualització de l'arbre de Huffman amb zoom i
          desplaçament. En mode per blocs es pot seleccionar quin bloc es
          vol inspeccionar.
    \item \textbf{Comparatives:} executa benchmarks sobre un directori de
          fitxers i mostra gràfiques de temps i taxa de compressió.
\end{itemize}

Durant la compressió o descompressió, la barra inferior mostra la fase
actual, el percentatge, els bytes processats i una estimació del temps
restant. Les operacions llargues s'executen amb \texttt{SwingWorker} per
no bloquejar la interfície.

\subsection{CLI}

La CLI exposa tres ordres:

\begin{itemize}
    \item \texttt{compress <input> <output>}
    \item \texttt{decompress <input> <output>}
    \item \texttt{benchmark <inputDir> <repetitions> <outputCsv>}
\end{itemize}

La tercera ordre és la que s'ha usat per a les mesures de la
Secció~\ref{sec:resultats}. La sortida CSV permet repetir proves, filtrar
resultats i actualitzar taules sense dependre d'una sessió manual de GUI.

\subsection{Ús Habitual}

El flux normal d'ús és el següent. L'usuari selecciona un fitxer des de
l'explorador integrat o escriu la ruta manualment. Si el fitxer no acaba
en \texttt{.hff}, la vista prepara automàticament una ruta de sortida
afegint aquesta extensió i el botó principal queda en mode compressió. Si
el fitxer d'entrada ja és un arxiu \texttt{.hff}, la vista canvia el botó
a mode extracció i proposa com a sortida el nom sense l'extensió.

Abans de comprimir, l'usuari pot triar la cua de prioritat i el mode de
compressió; per defecte s'utilitza heap binari i mode automàtic.
La resta d'opcions existeixen per experimentar i per poder
comparar estructures, no perquè sempre siguin millors.

Quan acaba l'operació, la pestanya de fitxers mostra les estadístiques
principals: mode efectiu, cua emprada, mida original, mida final,
metadades, càrrega, entropia, longitud mitjana i temps. Si el resultat és
per blocs, també mostra quants blocs han acabat en 1 byte, 2 bytes o
emmagatzemats. Les altres pestanyes permeten entrar en el detall: taula de
codis, arbre i blocs individuals.

%% ================================================================
\section{Disseny i Implementació}

\subsection{Flux de Compressió i Descompressió}

Els dos fluxos principals segueixen la mateixa idea: una seqüència de
passes petites on cada fase prepara la informació que necessita la
següent.

\textbf{Compressió}

\begin{enumerate}
    \item Comptar freqüències i construir plans candidats.
    \item Estimar la mida final de cada candidat, incloent capçalera,
          taula de freqüències i bits de càrrega.
    \item Seleccionar el mode demanat o, en \texttt{AUTO}, el més petit.
    \item Rebutjar la compressió si la mida estimada supera l'original.
    \item Escriure l'arxiu.
    \item Comprovar la mida real i eliminar la sortida si finalment és
          més gran que l'original.
\end{enumerate}

\textbf{Descompressió}

\begin{enumerate}
    \item Llegir la capçalera de l'arxiu \texttt{.hff}.
    \item Validar el format: identificador, mode, mida original i
          freqüències.
    \item Reconstruir la taula de freqüències i el mateix arbre de
          Huffman.
    \item Llegir les dades comprimides bit a bit i caminar per l'arbre
          fins arribar a fulles.
    \item Escriure els símbols originals fins arribar exactament a la
          mida original.
    \item Rebutjar l'arxiu si falta informació, sobren dades obligatòries
          o les mides reconstruïdes no quadren.
\end{enumerate}

La descompressió no reutilitza cap estat de la compressió anterior. Això
és important per al requisit de poder comprimir un dia i descomprimir un
altre: l'arxiu conté la capçalera, les freqüències i les metadades
necessàries per reconstruir l'arbre. Com que la construcció és
determinista, l'arbre reconstruït és equivalent al que es va usar durant
la compressió. En mode \texttt{HUFFMAN\_BLOCK}, el lector repeteix aquest
procés per a cada bloc: primer llegeix la mida i el mode local del bloc,
i després descodifica els bits d'aquell bloc amb l'arbre corresponent, o
copia les dades si el bloc és \texttt{STORED}. La mida original guardada
a la capçalera també indica quan s'ha d'aturar la lectura.

Els casos amb un únic símbol es tracten de forma especial. Si només hi ha
un símbol possible, no fa falta escriure cap bit de càrrega: la freqüència
ja diu quantes vegades s'ha de repetir aquell símbol.

\subsection{Modes de Compressió}

La Taula~\ref{tab:modes} resumeix els modes disponibles des del punt de
vista de disseny. La implementació concreta dels analitzadors es descriu
més endavant.

\begin{table}[H]
\centering
\caption{Modes de compressió implementats.}
\label{tab:modes}
\footnotesize
\begin{tabular}{@{}p{0.43\linewidth}p{0.47\linewidth}@{}}
\toprule
\textbf{Mode} & \textbf{Ús} \\
\midrule
\texttt{AUTO} & Analitza tots els modes Huffman i tria l'arxiu estimat més petit. \\
\texttt{HUFFMAN\_1\_BYTE} & Un arbre global sobre bytes individuals. \\
\texttt{HUFFMAN\_2\_BYTE} & Un arbre global sobre parelles de bytes; guarda el byte final separat si la mida és senar. \\
\texttt{HUFFMAN\_BLOCK} & Divideix el fitxer en blocs i tria localment entre 1 byte, 2 bytes i emmagatzemat. \\
\texttt{HUFFMAN\_BLOCK\_1\_BYTE} & Mode per blocs restringit a Huffman d'1 byte o emmagatzemat. \\
\texttt{HUFFMAN\_BLOCK\_2\_BYTE} & Mode per blocs restringit a Huffman de 2 bytes o emmagatzemat. \\
\texttt{STORED} & Només intern dins un bloc o per llegir arxius antics; no es tria per a fitxer sencer. \\
\bottomrule
\end{tabular}
\end{table}

El mode \texttt{AUTO} no prova a comprimir físicament el fitxer diverses
vegades. Primer construeix plans d'anàlisi. Un pla conté la taula de
freqüències, l'arbre, la taula de fulles per codificar ràpidament, el
nombre estimat de bits i la mida de les metadades. Amb això es pot
estimar la mida final abans d'escriure cap byte.

Per als modes globals, l'estimació és directa: mida de la capçalera més
bits de Huffman. Per al mode per blocs, la mida
estimada és la suma de la capçalera global i cada una de les millor opcions per a cada bloc.

Aquesta estratègia té dos avantatges. Primer, evita escriure arxius
temporals només per comparar modes. Segon, permet aplicar el requisit de
``no créixer'' abans de tocar el fitxer de sortida. Tot i així, després
d'escriure es fa una segona comprovació amb la mida real.

\subsection{Mode per Blocs}

El mode per blocs intenta resoldre el problema dels fitxers heterogenis:
parts del fitxer molt diferents que poden aprofitar de tenir un arbre
Huffman diferent. La decisió es pren localment, no amb una única taula de
freqüències per a tot el fitxer.

Per a cada bloc es comparen tres opcions lògiques:

\begin{itemize}
    \item guardar-lo sense comprimir (\texttt{STORED}),
    \item comprimir-lo amb Huffman d'1 byte,
    \item comprimir-lo amb Huffman de 2 bytes.
\end{itemize}

Una opció Huffman només s'accepta si guanya el cost d'emmagatzemar el
bloc directament. Això evita que la taula de freqüències faci créixer
blocs petits o poc compressibles.

\subsection{Format d'Arxiu}

L'arxiu propi comença amb una capçalera. Aquesta capçalera és el
motiu pel qual es pot tancar el programa després de comprimir i
descomprimir en una altra execució: l'arxiu conté la informació necessària
per reconstruir l'arbre.

\begin{table}[H]
\centering
\caption{Estructura simplificada de l'arxiu \texttt{.hff}.}
\label{tab:archive}
\begin{tabular}{@{}p{0.35\linewidth}p{0.55\linewidth}@{}}
\toprule
\textbf{Part} & \textbf{Contingut} \\
\midrule
Magic & Bytes \texttt{HUFF}, per detectar que el fitxer és del format esperat. \\
Mode & Identificador del mode escrit. \\
Mida original & Nombre de bytes que s'han de restaurar. També evita llegir bits extra com si fossin símbols reals. \\
Freqüències globals & Només en modes de fitxer sencer. Es guarden com a taula: símbol i freqüència. \\
Blocs & En mode per blocs, cada bloc guarda mida, mode, freqüències si cal i dades comprimides. \\
\bottomrule
\end{tabular}
\end{table}

\subsection{Recorregut de Classes}

Aquesta secció descriu les classes principals de la implementació. No
s'enumeren només per paquet: s'explica també com cooperen durant una
compressió, una descompressió i un benchmark.

L'organització general segueix el patró Model-Vista-Controlador, resumit a
la Fig.~\ref{fig:mvc} on es mostra l'organització general de paquets del projecte Java.
La vista concentra els components Swing i només envia
intencions; el controlador valida les peticions, crea el codec amb les
opcions triades i executa el treball en segon pla; el model conté la
compressió, el format d'arxiu, les cues, el progrés, la informació de
resultat i el benchmark. Aquesta separació evita que la GUI construeixi
arbres o interpreti arxius, i també permet que la CLI usi el mateix model
sense crear cap finestra.

\begin{figure*}[!b]
    \centering
    \resizebox{0.94\textwidth}{!}{%
    \begin{tikzpicture}[
        font=\sffamily\footnotesize,
        box/.style={draw, rounded corners, thick, align=center, minimum width=2.5cm, minimum height=0.9cm},
        model/.style={box, fill=blue!9},
        pkg/.style={model, minimum width=2.25cm},
        ui/.style={box, fill=green!10},
        ctrl/.style={box, fill=orange!15},
        tool/.style={box, fill=purple!9},
        core/.style={model, minimum width=3.0cm},
        arr/.style={->, thick}
    ]
        \node[ui] (view) at (-5.4, -0.45) {\texttt{view}\\GUI Swing};
        \node[tool] (cli) at (5.4, 2.2) {\texttt{CLI}\\línia de comandes};
        \node[ctrl] (controller) at (-2.4, 0.85) {\texttt{controller}\\\texttt{Controller}};
        \node[core] (codec) at (1.4, 0.85) {\texttt{HuffmanCodec}\\façana del model};
        \node[model] (benchmark) at (5.4, -0.45) {\texttt{benchmark}\\comparatives i CSV};

        \node[pkg] (codecPkg) at (-4.8, -1.8) {\texttt{codec}\\anàlisi i selecció};
        \node[pkg] (archive) at (-2.4, -1.8) {\texttt{archive}\\capçalera i modes};
        \node[pkg] (queue) at (0.0, -1.8) {\texttt{queue}\\cues de prioritat};
        \node[pkg] (progress) at (2.4, -1.8) {\texttt{progress}\\estat i fases};
        \node[pkg] (info) at (4.8, -1.8) {\texttt{info}\\dades per a la GUI};

        \draw[arr] (view) -- (controller);
        \draw[arr] (controller) -- (codec);
        \draw[arr] (cli) -- (codec);
        \draw[arr] (benchmark) -- (codec);
        \draw[arr] (codec) -- (codecPkg);
        \draw[arr] (codec) -- (archive);
        \draw[arr] (codec) -- (queue);
        \draw[arr] (codec) -- (progress);
        \draw[arr] (codec) -- (info);
        \draw[arr, dashed] (info.north) .. controls (2.9, -0.55) and (-2.9, -0.55) .. (view.east);
        \draw[arr, dashed] (progress) to[bend right=12] (controller);

        \node[ui, minimum width=1.8cm, minimum height=0.45cm, font=\sffamily\scriptsize] at (-5.15, 2.35) {Vista};
        \node[ctrl, minimum width=1.8cm, minimum height=0.45cm, font=\sffamily\scriptsize] at (-3.2, 2.35) {Controlador};
        \node[model, minimum width=1.8cm, minimum height=0.45cm, font=\sffamily\scriptsize] at (-1.05, 2.35) {Model};
        \node[tool, minimum width=1.8cm, minimum height=0.45cm, font=\sffamily\scriptsize] at (1.0, 2.35) {Eines};

    \end{tikzpicture}%
    }
    \caption{Arquitectura MVC per paquets i dependències principals.}
    \label{fig:mvc}
\end{figure*}

\subsubsection{Punts d'Entrada: Main i CLI}

\texttt{Main} és el punt d'entrada gràfic. La seva feina és mínima:
arrenca Swing al fil d'interfície, aplica l'aspecte visual del sistema,
crea una \texttt{MainWindow}, connecta un \texttt{Controller} i mostra la
finestra. No conté cap decisió de compressió. Aquesta simplicitat és
intencionada: si la classe d'entrada comença a crear arbres o llegir
fitxers, la separació MVC es perd des del principi.

\texttt{CLI} és l'adaptador textual. Implementa \texttt{run()} amb codi de
sortida i streams injectables, crea el mateix \texttt{HuffmanCodec} que la
GUI i només afegeix formatació de progrés i de resultats per terminal.

\subsubsection{Controller i ViewListener}

\texttt{ViewListener} és la interfície que la vista utilitza per comunicar
intencions: comprimir, descomprimir o executar comparatives. La vista no
coneix la classe \texttt{Controller}; només guarda una referència a aquesta
interfície.

\texttt{Controller} implementa \texttt{ViewListener}. En rebre una petició,
fa validacions bàsiques de ruta, neteja els resultats anteriors,
deshabilita controls i crea un \texttt{SwingWorker}. La part costosa viu a
\texttt{doInBackground()}, on es crea el \texttt{HuffmanCodec} amb la cua i
el mode triats. El codec rep un \texttt{ProgressListener}; cada vegada que
el model emet una mostra de progrés, el worker la publica. Després,
\texttt{process()} s'executa al fil de Swing i actualitza la barra de
progrés. Finalment, \texttt{done()} recupera el resultat i crida
\texttt{showCompressionInfo()}, \texttt{showDecompressionInfo()} o
\texttt{showBenchmarkReport()}.

Aquesta estructura evita dos problemes habituals. Primer, el compressor no
bloqueja la interfície. Segon, el model no actualitza components Swing
directament; només emet dades de progrés. El controlador fa de frontera
entre el treball de fons i la vista.

\subsubsection{Vista: MainWindow i Panells}

\texttt{MainWindow} és la finestra principal. Manté els camps de ruta, els
selectors de cua/mode i els panells descrits a la secció d'ús. Quan rep
un \texttt{CompressionInfo}, no torna a calcular res: simplement transforma
les dades en text, files de taula i nodes visuals.

La taula de codis usa un model propi, \texttt{SymbolTableModel}, amb
tipus de columna concrets. Això és important perquè les columnes de
freqüència i probabilitat s'ordenen numèricament, no com a text, i perquè
els símbols es mostren de forma llegible sense perdre el valor real que
necessita l'ordenació.

\texttt{FileBrowserPanel} implementa l'explorador integrat. Té un arbre de
directoris i una taula de fitxers. La taula usa comparadors propis per
ordenar nom, tipus, mida i data de modificació. En fer doble clic sobre
un directori, navega; en fer doble clic sobre un fitxer, avisa a
\texttt{MainWindow} per a que carregui la ruta seleccionada com a fitxer a processar.
També usa un \texttt{WatchService} per refrescar el directori actual quan canvia el
contingut, cosa útil després de crear un arxiu \texttt{.hff}.

\texttt{HuffmanTreePanel} dibuixa l'arbre. El model li passa un
\texttt{HuffmanTreeNodeInfo}; el panell calcula una disposició jeràrquica
recursiva i pinta arestes 0/1 i caixes de nodes. Manté escala i
desplaçament propis per permetre zoom amb roda de ratolí i arrossegament.
No coneix \texttt{HuffmanCode}; només coneix la representació d'informació
preparada pel model.

\texttt{BlockListPanel} mostra una llista lateral de blocs en mode
\texttt{HUFFMAN\_BLOCK}. Cada cel·la indica número, mode local i mida, amb
colors diferents per 1 byte, 2 bytes i emmagatzemat. Quan l'usuari tria
un bloc, la finestra rep el \texttt{BlockInfo} corresponent i actualitza
la taula de codis i l'arbre amb les dades d'aquell bloc.

\texttt{BenchmarkPanel} conté el formulari de directori/repeticions, la barra
de progrés, el botó d'exportació CSV i les pestanyes de gràfiques.
\texttt{BenchmarkGraphPanel} és un panell de dibuix 2D simple: rep sèries
de punts, calcula límits, genera marques d'eix amb \texttt{ChartAxisSupport}
i pinta línies, punts, etiquetes i llegenda.

\subsubsection{Format: CompressionMode i ArchiveHeader}

\texttt{CompressionMode} concentra tots els noms de mode en un únic
\texttt{enum}. No són tipus Java separats, però no tots els valors són
vàlids en tots els punts del programa. Per això la classe defineix
mètodes de context, com \texttt{requestModes()}, \texttt{isArchiveMode()}
i \texttt{allowedBlockHuffmanModes()}, que indiquen on es pot usar cada
valor.

Els modes que es poden escriure a la capçalera \texttt{.hff} són
\texttt{STORED}, \texttt{HUFFMAN\_1\_BYTE}, \texttt{HUFFMAN\_2\_BYTE} i
\texttt{HUFFMAN\_BLOCK}. Aquests són els únics amb identificador binari,
perquè són els que el lector ha de trobar quan obre un arxiu. En canvi,
\texttt{AUTO}, \texttt{HUFFMAN\_BLOCK\_1\_BYTE} i
\texttt{HUFFMAN\_BLOCK\_2\_BYTE} són ordres de compressió: serveixen per
dir al selector què ha d'analitzar, però no descriuen un format nou dins
el fitxer.

Aquesta distinció és necessària encara que tot visqui dins la mateixa
classe. \texttt{AUTO} pot aparèixer a la GUI perquè és una opció útil per
a l'usuari, però no pot arribar al writer com a mode de capçalera perquè
un lector futur no sabria quin arbre reconstruir. Els submodes
\texttt{HUFFMAN\_BLOCK\_1\_BYTE} i
\texttt{HUFFMAN\_BLOCK\_2\_BYTE} només restringeixen l'anàlisi dels blocs;
si guanyen, l'arxiu continua guardant \texttt{HUFFMAN\_BLOCK}. Finalment,
\texttt{STORED} és serialitzable perquè pot aparèixer com a bloc local
emmagatzemat o en arxius antics, però no s'ofereix com a compressió de
fitxer sencer perquè sempre afegiria capçalera i faria créixer el
resultat.

\texttt{ArchiveHeader} escriu i llegeix la capçalera binària. Guarda
quatre bytes que identifiquen la capçalera (\textit{magic}), el mode, la mida original i, si el mode és global,
una taula de freqüències. No guarda sempre les 256 o 65.536 freqüències;
escriu només les entrades presents (tenen freqüència major a 0). En llegir, valida la
capçalera abans de deixar continuar: \textit{magic} incorrecte, mode desconegut,
freqüències duplicades o totals incoherents generen
\texttt{ArchiveFormatException}.

\subsubsection{FrequencyTable i HuffmanCode}

\texttt{FrequencyTable} és la representació bàsica de les freqüències.
Internament usa un array \texttt{long[]} de mida fixa: 256 posicions per
bytes o 65.536 per parelles. Aquesta decisió consumeix més memòria que un
mapa en alguns casos, però dóna accés directe $O(1)$ i fa molt
simple calcular mides, entropia i taules de capçalera. A més guarda
\texttt{totalCount} i \texttt{distinctSymbolCount}; així no cal recórrer
tot l'array cada vegada que es vol saber quants símbols diferents hi ha.

El mètode \texttt{addSymbol()} incrementa la freqüència del símbol. Si la
freqüència anterior era zero, també incrementa \texttt{distinctSymbolCount}.
Aquesta comprovació serveix per mantenir actualitzat el nombre de símbols
diferents sense haver de recórrer tot l'array després de cada bloc o
estimació. Aquest comptador s'utilitza per saber si hi ha el cas especial
d'un únic símbol i per calcular ràpidament el cost de guardar la taula de
freqüències a la capçalera. \texttt{entropy()} calcula l'entropia de la
distribució observada. \texttt{encodedBitCount()} multiplica la freqüència
de cada símbol per la profunditat de la seva fulla de Huffman i obté el
nombre exacte de bits de càrrega. \texttt{clear()} permet reutilitzar la
mateixa taula en la selecció de mida de bloc, evitant crear repetidament
arrays de 65.536 posicions.

\texttt{HuffmanCode} representa tant fulles com nodes interns de l'arbre.
Una fulla conté un símbol real i una freqüència. Un node intern conté els
dos fills i la suma de freqüències. Quan dos nodes tenen la mateixa
freqüència, \texttt{compareTo()} desempata pel símbol mínim del subarbre.
Amb això, l'arbre és determinista independentment de l'implementació de
la cua de prioritat.

La construcció de l'arbre es fa a \texttt{buildFromFrequencies()}. El
mètode rep una \texttt{FrequencyTable} i una \texttt{NodeQueue} que
és l'interfície comuna de les cues de prioritat del
projecte: l'algorisme només necessita inserir nodes, extreure el mínim i
consultar la mida, independentment de si al darrere hi ha un heap binari,
una llista ordenada o un heap de Fibonacci.

Els mètodes \texttt{single()} i \texttt{merge()} materialitzen els dos casos
del model teòric: fulles amb símbol real i nodes interns amb dos fills. La
part rellevant aquí no és repetir l'algorisme, sinó que tots els nodes
passen per \texttt{compareTo()}, de manera que el desempat determinista és
compartit per totes les cues.

El codi binari de cada fulla no es construeix durant aquest procés. Es
calcula després de forma mandrosa amb una única passada en amplada: cada
fill afegeix un 0 o un 1 al codi del pare i guarda també la profunditat.
Finalment, \texttt{lookupBySymbol()} crea una taula
símbol $\rightarrow$ fulla perquè el writer pugui codificar bytes sense
recórrer l'arbre per cada símbol.


\subsubsection{Cues de Prioritat}

L'algorisme de Huffman només necessita tres operacions de les cues: inserir,
extreure el mínim i consultar la mida. Per això el model defineix una
interfície mínima \texttt{NodeQueue}. El codi que construeix l'arbre no
sap quina estructura concreta hi ha darrere d'aquesta interfície; només
usa aquestes tres operacions.

\texttt{PriorityQueueStrategy} és l'\texttt{enum} que connecta la selecció
de l'usuari amb la implementació real. La GUI i el benchmark treballen amb
aquests valors, com \texttt{BINARY\_HEAP}, \texttt{DICHOTOMIC\_LIST} i
\texttt{FIBONACCI\_HEAP}. Quan el codec necessita construir un arbre, la
estratègia crea la \texttt{NodeQueue} corresponent amb \texttt{createQueue()}.
Així el selector visible a la interfície, el benchmark i el codi de
Huffman comparteixen una sola llista de cues disponibles.

Les tres implementacions són:

\begin{itemize}
    \item \textbf{\texttt{BinaryHeapNodeQueue}:} estratègia per defecte,
          basada en una \texttt{PriorityQueue} de Java. Internament aquesta
          estructura manté els elements en un heap binari: el mínim queda
          sempre a l'arrel i, quan s'insereix o s'elimina, es reordena
          només el camí necessari per conservar la propietat de heap. Té
          bon rendiment pràctic, poca lògica pròpia i cost $O(\log n)$ per
          inserció i extracció.
    \item \textbf{\texttt{DichotomicListNodeQueue}:} manté una llista
          ordenada amb \texttt{ArrayList}. Quan s'afegeix un node, usa
          cerca binària per trobar la posició que li correspon segons
          l'ordre de Huffman i l'insereix en aquell punt. Això fa que
          extreure el mínim sigui directe, perquè el mínim és al principi
          de la llista. El cost amagat és el desplaçament dels elements de
          l'array: trobar la posició és ràpid, però inserir pot costar
          $O(n)$.
    \item \textbf{\texttt{FibonacciHeapNodeQueue}:} estructura més
          complexa, amb insercions molt barates i extracció mínima
          amortitzada logarítmica. La implementació
          manté una llista circular d'arrels i cada node pot tenir una
          llista circular de fills. Inserir és gairebé només afegir una
          arrel nova i actualitzar el punter al mínim. En extreure el
          mínim, els seus fills passen a la llista d'arrels i després es
          fa una consolidació: si dues arrels tenen el mateix grau,
          s'uneixen i la de menor prioritat queda com a pare.
\end{itemize}


\subsubsection{Plans Interns i HuffmanArchive}

\texttt{CodecPlans} no conté lògica; conté tres \textit{records} interns que fan
de frontera entre anàlisi i escriptura.

\texttt{WholePlan} descriu un
mode global complet: freqüències, arrel, fulles, bits estimats, entropia,
longitud mitjana, capçalera i mida estimada. S'usa tant per al mode
d'1 byte com per al de 2 bytes.

\texttt{BlockUnit} descriu un bloc concret dins el mode per blocs. Guarda
la mida del bloc, el mode local seleccionat i les dades que el writer
necessitarà. Pot contenir la taula/arbre d'1 byte, la taula/arbre de
2 bytes o cap arbre si el bloc ha quedat en mode \texttt{STORED} emmagatzemat.
La raó és que l'analitzador ja ha hagut de construir aquests arbres per saber quants
bits ocuparia cada alternativa i decidir el mode local del bloc. Si el
writer només guardàs el mode i les freqüències, hauria de tornar a
construir el mateix arbre durant l'escriptura. Això duplicaria temps de
càlcul i faria que l'escriptura depengués d'una segona reconstrucció que
no aporta cap decisió nova. Amb \texttt{BlockUnit}, el writer rep el pla
tancat: només serialitza exactament el que l'anàlisi ha triat.

\texttt{BlockPlan} agrupa tots els \texttt{BlockUnit} i afegeix dades
agregades del mode per blocs: mida de bloc triada, capçalera, metadades,
mida estimada, bits totals Huffman, entropia mitjana ponderada i longitud
mitjana.

\texttt{HuffmanArchive} representa l'arxiu Huffman en el punt on ja se
sap quin format s'ha triat. Durant la compressió, els analitzadors poden
generar diversos plans candidats, però només un s'acabarà escrivint.
\texttt{HuffmanArchive} empaqueta aquesta decisió: guarda el mode final,
la capçalera i el pla concret que correspon a aquell mode. Si el mode és
global, conté un \texttt{WholePlan}; si és per blocs, conté un
\texttt{BlockPlan}. Això evita passar per separat un mode, una capçalera i
un pla opcional a cada mètode del writer.

La mateixa classe també s'utilitza després de llegir un arxiu. En aquest
cas no conté un pla de compressió, perquè la decisió ja havia estat presa
quan es va crear el fitxer. Conté la capçalera llegida i la informació
reconstruïda que necessita la GUI: símbols, arbre i blocs. Per això
\texttt{HuffmanArchive} és una frontera comuna entre el codec i el
format: abans d'escriure diu ``això és el que s'ha de serialitzar'', i
després de llegir diu ``això és el que s'ha reconstruït del fitxer''.

\subsubsection{Analitzadors i Selector}

\texttt{HuffmanCodec} és una façana. Manté l'API pública que usen la GUI,
la CLI i el benchmark, però delega les parts concretes:
\texttt{HuffmanWholeFileAnalyzer} analitza modes globals,
\texttt{HuffmanBlockAnalyzer} analitza blocs, \texttt{HuffmanModeSelector}
tria el pla final, i \texttt{HuffmanArchiveWriter}/\texttt{Reader}
només escriuen o llegeixen el format.

La raó d'aquesta divisió és evitar que la classe principal del codec
barregi tres preguntes diferents: ``quant ocuparia cada opció?'', ``quina
opció convé escriure?'' i ``com s'escriu exactament aquesta opció?''. Si
aquestes decisions viuen juntes, el codi sembla més curt al principi però
és més difícil d'explicar i provar. Separant-les, cada classe té una
responsabilitat que coincideix amb una part del problema.

\texttt{HuffmanAnalyzer} és una superclasse petita per als analitzadors.
Només centralitza \texttt{buildTree()}: construeix l'arbre amb la cua triada
i acumula el temps invertit en aquesta construcció per tal de poder comparar
el seu rendiment. No força cap flux
d'anàlisi comú, perquè analitzar un fitxer sencer i analitzar blocs són
processos diferents.

\texttt{HuffmanWholeFileAnalyzer} construeix els plans globals.
\texttt{analyzeByte()} llegeix el fitxer en buffers i afegeix cada byte a
una \texttt{FrequencyTable} de 256 símbols. \texttt{analyzeWord()} fa el
mateix amb parelles de bytes. En aquest cas, ``buffer'' vol dir buffer
d'entrada de Java, no bloc de compressió: és només la manera de llegir el
fitxer a trossos. Si una lectura acaba just després del primer byte d'una
parella, aquest byte no es pot comptar encara com a símbol de 16 bits; es
guarda a \texttt{pendingByte}. En llegir el buffer següent, el primer byte
nou es combina amb aquest pendent i la parella ja es pot afegir a la
taula de freqüències. Si el fitxer acaba i encara queda un
\texttt{pendingByte}, vol dir que la mida original és senar: aquest byte
final no forma cap parella i es guarda separat al pla perquè el writer
l'escrigui cru al final del flux Huffman. Després de comptar, construeix
l'arbre, calcula fulles, bits, entropia, capçalera i mida estimada.

\texttt{HuffmanBlockSizeSelector} tria una mida de bloc abans de fer
l'anàlisi complet. Les mides candidates són 4 KiB, 16 KiB, 64 KiB,
256 KiB, 1 MiB, 4 MiB i 16 MiB; totes són múltiples de 4 KiB. Per això el
selector pot llegir el fitxer en mostres de 4 KiB i alimentar totes les
candidates al mateix temps: la candidata de 4 KiB tanca un bloc a cada
mostra, la de 16 KiB cada quatre mostres, la de 64 KiB cada setze, i així
successivament. Cada candidata acumula freqüències fins omplir el seu
bloc sintètic; llavors estima si convindria emmagatzemar, usar Huffman
d'1 byte o usar Huffman de 2 bytes. L'estimació usa entropia i cost de
metadades, no arbres reals, perquè construir arbres per a totes les
candidates seria massa car.

\texttt{HuffmanBlockAnalyzer} fa l'anàlisi complet del mode per blocs amb
la mida ja seleccionada. Recorre el fitxer bloc a bloc i, per a cadascun,
prova només els candidats Huffman que permet el mode demanat:
\texttt{HUFFMAN\_BLOCK} i \texttt{AUTO} permeten 1 byte i 2 bytes,
\texttt{HUFFMAN\_BLOCK\_1\_BYTE} només 1 byte, i
\texttt{HUFFMAN\_BLOCK\_2\_BYTE} només 2 bytes. \texttt{STORED} sempre
queda disponible com a alternativa local.

Per cada candidat permès, l'analitzador calcula el cost del bloc:
capçalera local, taula de freqüències i bits Huffman. En el candidat de
2 bytes també suma el byte final cru si el bloc té mida senar. Després
compara aquests costos amb el cost de copiar el bloc directament. Si cap
candidat Huffman és més petit que copiar-lo, el bloc es marca com a
\texttt{STORED}. Cada bloc és independent: no es comparteixen parelles de
2 bytes ni arbres entre blocs.

\texttt{HuffmanModeSelector} és el punt on es converteix l'anàlisi en una
decisió final. Rep plans ja estimats i tria el més petit en mode
\texttt{AUTO}, o bé comprova que el mode demanat per l'usuari sigui
aplicable. Així la lògica de selecció queda separada de la lectura del
fitxer i de l'escriptura del format.

\subsubsection{Writer i Reader}

\texttt{HuffmanArchiveWriter} rep un \texttt{HuffmanArchive} ja triat i
escriu exactament el pla seleccionat.

\texttt{BitOutputStream} és el suport per escriure bits. Acumula bits en
un byte temporal i el descarrega quan està ple. Si l'últim byte queda
incomplet, l'omple amb zeros. Aquests zeros són segurs perquè el decoder
sap quants bytes originals ha de restaurar i s'atura quan arriba a aquesta
mida.

\texttt{HuffmanArchiveReader} fa el camí invers. Llegeix
\texttt{ArchiveHeader}, reconstrueix la taula de freqüències i torna a
construir l'arbre. En modes Huffman, usa \texttt{BitInputStream} per
llegir un bit cada vegada i caminar per l'arbre fins trobar una fulla. En
mode per blocs, repeteix aquest procés per cada bloc i valida que la suma
de bytes restaurats arribi exactament a la mida original.

\subsubsection{Informació, Progrés i Utilitats}

El paquet \texttt{info} conté dades immutables. \texttt{CompressionStats}
i \texttt{DecompressionStats} són el resum numèric. \texttt{CompressionInfo}
i \texttt{DecompressionInfo} combinen estadístiques amb taula de símbols,
arbre i llista de blocs. \texttt{HuffmanSymbolInfo},
\texttt{HuffmanTreeNodeInfo} i \texttt{BlockInfo} són còpies preparades
per mostrar: contenen només els camps que necessita la GUI, com símbol,
freqüència, probabilitat, codi, fills de l'arbre, mida de bloc o mode
local. La vista no rep les estructures internes que usa el codec per
comprimir, sinó aquests objectes de només lectura. \texttt{HuffmanInfoFactory}
fa aquesta conversió, calculant codis com a cadenes i probabilitats per a
la taula.

El paquet \texttt{progress} desacobla el model de la interfície.
\texttt{ProgressPhase} enumera les fases visibles del treball, com
analitzar, escriure o llegir. \texttt{ProgressTracker} rep comptadors de
bytes i emet \texttt{ProgressSnapshot} cada 64 KiB o quan una fase acaba.
L'estimació del temps restant és una extrapolació simple: temps
transcorregut, bytes processats i bytes pendents. Si encara no hi ha prou
informació, retorna \texttt{-1} i la vista mostra ``--''.

\texttt{ByteFormat} només converteix mides de bytes a unitats llegibles
per la GUI i la CLI.

\subsubsection{Benchmark}

\texttt{BenchmarkConfig} guarda la configuració d'una execució: directori
d'entrada, nombre de repeticions i opcions seleccionades.
\texttt{BenchmarkService} carrega els fitxers del directori, executa les
comparatives descrites a la secció de resultats i verifica cada cicle de
compressió i descompressió
amb \texttt{Files.mismatch()}. Durant l'execució,
\texttt{BenchmarkProgressListener} rep mostres de progrés perquè la GUI o
la CLI puguin mostrar l'estat sense que el servei conegui cap component
visual. Finalment, \texttt{BenchmarkReport} agrupa les files i
\texttt{BenchmarkCsvWriter} les serialitza en CSV.


%% ================================================================
\section{Validació amb Proves Unitàries}

Les proves s'executen amb JUnit 5 mitjançant \texttt{mvn test}. En total hi
ha 31 proves i no creen ni llegeixen cap corpus. La
validació s'ha organitzat al voltant de tres riscos principals: perdre
dades en comprimir i descomprimir, construir arbres diferents amb cues diferents i
acceptar arxius corruptes com si fossin vàlids.

El grup més gran és \texttt{HuffmanCodecTest}, amb 20 proves. Aquestes
proves exerciten el compressor complet: escriuen dades originals, les
comprimeixen, les descomprimeixen i comparen el resultat byte a byte.
Inclouen fitxers d'un sol símbol, els 256 valors possibles de byte, text
real, símbols de 2 bytes i modes per blocs. També comproven que el mode
automàtic rebutja fitxers buits o casos on el resultat creixeria, que els
avisos de progrés avancen de forma consistent i que l'exportador CSV del
benchmark manté les capçaleres i files esperades.

La robustesa del format es valida amb casos negatius dins el mateix grup de
codec. Es creen arxius amb capçalera incorrecta, modes no suportats,
càrrega truncada o totals de freqüència incompatibles amb la mida original.
En tots aquests casos s'espera una excepció de format, no una sortida
parcial. Aquesta part és important perquè un compressor pot fallar de forma
silenciosa si accepta bits insuficients o metadades incoherents.

\texttt{NodeQueueTest} conté 9 proves independents del compressor. Cada cua
ha d'eliminar nodes en ordre creixent, comportar-se igual que
\texttt{PriorityQueue} sota 2.000 operacions mixtes aleatòries i fallar
correctament quan es demana extreure el mínim d'una cua buida. Les mateixes
proves s'apliquen al heap binari, la llista dicotòmica i el heap de
Fibonacci, perquè totes tres han de respectar el mateix contracte.

Finalment, \texttt{HuffmanCodeTest} cobreix dues propietats internes de
l'arbre. La primera és que, després de construir-lo, cada símbol existent es
pot localitzar directament a partir del seu valor i conserva la freqüència
esperada. La segona força el càlcul posterior dels codis i comprova que la
profunditat de cada fulla coincideix amb la longitud del codi binari
assignat.

%% ================================================================
\section{Resultats i Discussió}
\label{sec:resultats}

\subsection{Protocol de Benchmark}

Per a aquesta primera versió s'ha utilitzat el corpus Silesia, un conjunt de
12 fitxers pensat per provar
compressors sense pèrdues amb dades de naturalesa diferent: text, HTML,
executables, bases de dades i imatges mèdiques. La mida total original és
d'aproximadament 212 MB \cite{ref_silesia}.

La mesura s'ha repetit tres vegades per a cada combinació de fitxer i
configuració. El CSV exportat conté la mitjana d'aquestes tres repeticions;
a partir d'aquestes files, la memòria presenta també medianes entre els 12
fitxers per reduir l'efecte dels fitxers molt grans. La comanda equivalent
és:

\begin{lstlisting}[language=bash,caption={Benchmark CLI sobre Silesia.},label={lst:bench}]
java -cp target/P4-1.0-SNAPSHOT.jar \
  com.serafinebot.p4.CLI benchmark \
  <silesiaDir> 3 memoria/silesia-benchmark.csv
\end{lstlisting}

La GUI usa exactament el mateix \texttt{BenchmarkService}. La diferència
és que la CLI exporta directament el CSV i evita interferència manual. El
directori emprat conté només els 12 fitxers de Silesia, sense fitxers ocults
del sistema.

El benchmark fa dues comparatives:

\begin{itemize}
    \item \textbf{Cues de prioritat:} executa \texttt{AUTO} amb cada
          estructura i mesura el temps de construcció d'arbre en compressió
          i descompressió.
    \item \textbf{Modes de compressió:} amb heap binari, compara els
          modes Huffman d'1 byte, 2 bytes i blocs restringits.
\end{itemize}

Cada compressió es descomprimeix després i es compara amb l'original.
Per tant, una fila de benchmark només és vàlida si també ha passat la
verificació de descompressió.

\subsection{Resultats per Cua de Prioritat}

La Taula~\ref{tab:queueavg} compara només el temps de construir arbres en les fases
de compressió i descompressió.
La reducció mitjana és idèntica perquè el desempat determinista fa que les
tres cues produeixin els mateixos arbres finals. També aquí es mostren
mitjanes i medianes: la mediana és útil perquè alguns fitxers tenen molt
pocs símbols diferents i fan que el temps de construcció sigui quasi zero.

\begin{table}[!b]
\centering
\caption{Temps de construcció d'arbre per cua.}
\label{tab:queueavg}
\scriptsize
\begin{tabular}{@{}lrrrrr@{}}
\toprule
\textbf{Cua} & \textbf{C. mitj.} & \textbf{C. med.} & \textbf{D. mitj.} & \textbf{D. med.} & \textbf{Red.} \\
\midrule
Binari & 26,833 & 6,333 & 5,556 & 1,000 & 40,315 \\
Fibonacci & 133,444 & 40,167 & 29,139 & 5,167 & 40,315 \\
Dicotòmica & 450,000 & 28,500 & 99,500 & 1,500 & 40,315 \\
\bottomrule
\end{tabular}
\end{table}

\begin{figure}[t]
    \centering
    \pgfplotstableread[col sep=comma]{
queue,compmean,compmedian,decompmean,decompmedian
Binari,26.833,6.333,5.556,1.000
Fibonacci,133.444,40.167,29.139,5.167
Llista,450.000,28.500,99.500,1.500
    }\queuestatsdata
    \begin{tikzpicture}
    \begin{axis}[
        ybar,
        width=\linewidth,
        height=5.2cm,
        ylabel={Temps (ms)},
        symbolic x coords={Binari,Fibonacci,Llista},
        xtick=data,
        x tick label style={font=\scriptsize},
        ymode=log,
        log basis y={10},
        bar width=4pt,
        enlarge x limits=0.18,
        legend columns=2,
        legend pos=north west,
        legend style={font=\tiny, fill=white, fill opacity=0.85, draw opacity=1},
        tick label style={font=\scriptsize},
        label style={font=\scriptsize},
        grid=major
    ]
        \addplot table[x=queue,y=compmean] {\queuestatsdata};
        \addlegendentry{Comp. mitjana}
        \addplot table[x=queue,y=compmedian] {\queuestatsdata};
        \addlegendentry{Comp. mediana}
        \addplot table[x=queue,y=decompmean] {\queuestatsdata};
        \addlegendentry{Descomp. mitjana}
        \addplot table[x=queue,y=decompmedian] {\queuestatsdata};
        \addlegendentry{Descomp. mediana}
    \end{axis}
    \end{tikzpicture}
    \caption{Diferència entre mitjana i mediana en la construcció d'arbres durant la compressió i la descompressió.}
    \label{fig:queue_mean_median}
\end{figure}

El heap binari és clarament la millor opció pràctica per defecte. A més del
seu cost asimptòtic raonable, és la implementació interna de Java
(\texttt{PriorityQueue}), molt optimitzada i amb poca sobrecàrrega constant.
En canvi, el heap de Fibonacci i la llista dicotòmica són implementacions
pròpies del projecte i, per tant, no estan tant optimitzades.
El heap de Fibonacci té bones propietats teòriques, però en aquesta aplicació la
complexitat de l'estructura i la manipulació d'enllaços pesen més que el
guany amortitzat. La Fig.~\ref{fig:queue_mean_median} també mostra que la
mediana és menor que la mitjana en totes les cues, perquè alguns fitxers
tenen pocs símbols diferents i fan poca feina de construcció. La llista
dicotòmica és la més perjudicada en la mitjana perquè, quan el nombre de
símbols creix, mantenir la llista ordenada implica desplaçar més elements.

\subsection{Resultats per Mode}

La Taula~\ref{tab:bestmode} mostra el mode que més comprimeix cada fitxer.
El percentatge indica reducció respecte a l'original: per exemple, 47\%
significa que l'arxiu final ocupa un 47\% menys. Els temps són la mitjana
de tres execucions d'aquell mode sobre el mateix fitxer.

La Fig.~\ref{fig:mode_ratio_by_file} mostra que el mode de 2 bytes és molt
competitiu en aquests fitxers. Això indica que molts fitxers tenen
repeticions útils a nivell de parelles de bytes. El mode per blocs guanya
especialment quan la distribució local varia: \texttt{xml}, \texttt{reymont},
\texttt{x-ray}, \texttt{mr} i \texttt{samba} milloren amb blocs de 2 bytes,
mentre que \texttt{ooffice} prefereix blocs d'1 byte. També s'observa que
els modes de 2 bytes no sempre superen els blocs: en fitxers heterogenis,
tenir arbres locals pot compensar la capçalera extra de cada bloc.

\begin{figure}[H]
    \centering
    \pgfplotstableread[col sep=comma]{
mode,ratio,comp,decomp
1 byte,31.261,696.2,326.8
2 bytes,39.340,751.9,413.5
Bloc 1 byte,33.346,471.8,323.7
Bloc 2 bytes,39.866,710.1,404.4
    }\modedata
    \begin{tikzpicture}
    \begin{axis}[
        ybar,
        width=\linewidth,
        height=5.3cm,
        ylabel={Reducció mitjana (\%)},
        symbolic x coords={1 byte,2 bytes,Bloc 1 byte,Bloc 2 bytes},
        xtick=data,
        x tick label style={rotate=20,anchor=east},
        ymin=0,
        ymax=45,
        bar width=9pt,
        grid=major
    ]
        \addplot table[x=mode,y=ratio] {\modedata};
    \end{axis}
    \end{tikzpicture}
    \caption{Reducció mitjana per mode de compressió sobre Silesia.}
    \label{fig:mode_ratio}
\end{figure}

\begin{figure*}[p]
\centering
\footnotesize

\refstepcounter{table}\label{tab:bestmode}
\textbf{Taula~\thetable.} Millor mode per fitxer en el benchmark Silesia.\par\smallskip
\scriptsize
\begin{tabular}{@{}lrlrr@{}}
\toprule
\textbf{Fitxer} & \textbf{Mida (B)} & \textbf{Millor mode} & \textbf{Reducció (\%)} & \textbf{Comp. / Descomp. (ms)} \\
\midrule
xml & 5.345.280 & \texttt{HUFFMAN\_BLOCK\_2\_BYTE} & 47,051 & 168,6 / 94,9 \\
ooffice & 6.152.192 & \texttt{HUFFMAN\_BLOCK\_1\_BYTE} & 18,712 & 221,3 / 169,6 \\
reymont & 6.627.202 & \texttt{HUFFMAN\_BLOCK\_2\_BYTE} & 49,824 & 207,3 / 111,4 \\
sao & 7.251.944 & \texttt{HUFFMAN\_2\_BYTE} & 12,316 & 468,1 / 353,1 \\
x-ray & 8.474.240 & \texttt{HUFFMAN\_BLOCK\_2\_BYTE} & 29,840 & 366,8 / 275,1 \\
mr & 9.970.564 & \texttt{HUFFMAN\_BLOCK\_2\_BYTE} & 61,019 & 331,9 / 151,4 \\
osdb & 10.085.684 & \texttt{HUFFMAN\_2\_BYTE} & 21,563 & 597,6 / 416,7 \\
dickens & 10.192.446 & \texttt{HUFFMAN\_2\_BYTE} & 49,375 & 342,6 / 223,3 \\
samba & 21.606.400 & \texttt{HUFFMAN\_BLOCK\_2\_BYTE} & 35,748 & 950,4 / 454,9 \\
nci & 33.553.445 & \texttt{HUFFMAN\_2\_BYTE} & 72,365 & 961,6 / 275,4 \\
webster & 41.458.703 & \texttt{HUFFMAN\_2\_BYTE} & 47,301 & 1374,7 / 830,8 \\
mozilla & 51.220.480 & \texttt{HUFFMAN\_2\_BYTE} & 35,224 & 2533,6 / 1434,1 \\
\bottomrule
\end{tabular}
\par\medskip

\refstepcounter{table}\label{tab:modeavg}
\textbf{Taula~\thetable.} Mitjanes i medianes per mode sobre els 12 fitxers.\par\smallskip
\scriptsize
\begin{tabular}{@{}lrrrrrr@{}}
\toprule
\textbf{Mode} & \textbf{Red. mitj.} & \textbf{Red. med.} & \textbf{Comp. mitj.} & \textbf{Comp. med.} & \textbf{Descomp. mitj.} & \textbf{Descomp. med.} \\
\midrule
\texttt{HUFFMAN\_1\_BYTE} & 31,261 & 26,981 & 696,2 & 404,4 & 326,8 & 244,8 \\
\texttt{HUFFMAN\_2\_BYTE} & 39,340 & 38,982 & 751,9 & 472,5 & 413,5 & 273,7 \\
\texttt{BLOCK\_1\_BYTE} & 33,346 & 34,665 & 471,8 & 254,6 & 323,7 & 248,3 \\
\texttt{BLOCK\_2\_BYTE} & 39,866 & 41,399 & 710,1 & 425,7 & 404,4 & 276,3 \\
\bottomrule
\end{tabular}
\par\medskip

\refstepcounter{figure}\label{fig:mode_ratio_by_file}
\textbf{Fig.~\thefigure.} Reducció obtinguda per cada mode en cada fitxer de Silesia.\par\smallskip
    \begin{tikzpicture}
    \begin{axis}[
        width=0.98\textwidth,
        height=6cm,
        ylabel={Reducció (\%)},
        xlabel={Mida del fitxer (MB)},
        xmin=0,
        xmax=54,
        ymin=0,
        ymax=80,
        legend columns=4,
        legend style={font=\footnotesize, at={(0.5,1.03)}, anchor=south, draw=none},
        grid=major
    ]
        \addplot+[mark=*] coordinates {(5.35,30.555) (6.15,16.671) (6.63,39.135) (7.25,5.603) (8.47,17.118) (9.97,53.611) (10.09,17.267) (10.19,42.832) (21.61,23.408) (33.55,69.528) (41.46,37.456) (51.22,21.947)};
        \addlegendentry{1 byte}
        \addplot+[mark=square*] coordinates {(5.35,42.740) (6.15,18.622) (6.63,49.785) (7.25,12.316) (8.47,29.386) (9.97,60.978) (10.09,21.563) (10.19,49.375) (21.61,32.429) (33.55,72.365) (41.46,47.301) (51.22,35.224)};
        \addlegendentry{2 bytes}
        \addplot+[mark=triangle*] coordinates {(5.35,36.914) (6.15,18.712) (6.63,39.469) (7.25,5.950) (8.47,19.210) (9.97,53.700) (10.09,17.267) (10.19,42.874) (21.61,32.416) (33.55,69.527) (41.46,37.542) (51.22,26.566)};
        \addlegendentry{Bloc 1 byte}
        \addplot+[mark=diamond*] coordinates {(5.35,47.051) (6.15,18.622) (6.63,49.824) (7.25,12.316) (8.47,29.840) (9.97,61.019) (10.09,21.563) (10.19,49.375) (21.61,35.748) (33.55,72.346) (41.46,47.244) (51.22,33.448)};
        \addlegendentry{Bloc 2 bytes}
        \node[font=\tiny, rotate=45, anchor=west] at (axis cs:5.35,47.051) {xml};
        \node[font=\tiny, rotate=45, anchor=west] at (axis cs:6.15,18.622) {ooffice};
        \node[font=\tiny, rotate=45, anchor=west] at (axis cs:6.63,49.824) {reymont};
        \node[font=\tiny, rotate=45, anchor=west] at (axis cs:7.25,12.316) {sao};
        \node[font=\tiny, rotate=45, anchor=west] at (axis cs:8.47,29.840) {x-ray};
        \node[font=\tiny, rotate=45, anchor=west] at (axis cs:9.97,61.019) {mr};
        \node[font=\tiny, rotate=45, anchor=west] at (axis cs:10.09,21.563) {osdb};
        \node[font=\tiny, rotate=45, anchor=west] at (axis cs:10.19,49.375) {dickens};
        \node[font=\tiny, rotate=45, anchor=west] at (axis cs:21.61,35.748) {samba};
        \node[font=\tiny, rotate=45, anchor=west] at (axis cs:33.55,72.346) {nci};
        \node[font=\tiny, rotate=45, anchor=west] at (axis cs:41.46,47.244) {webster};
        \node[font=\tiny, rotate=45, anchor=west] at (axis cs:51.22,33.448) {mozilla};
    \end{axis}
    \end{tikzpicture}
\par\medskip

\refstepcounter{figure}\label{fig:mode_time_by_file}
\textbf{Fig.~\thefigure.} Temps mitjà de compressió per mode en cada fitxer de Silesia.\par\smallskip
    \begin{tikzpicture}
    \begin{axis}[
        width=0.98\textwidth,
        height=6cm,
        ylabel={Temps de compressió (ms)},
        xlabel={Mida del fitxer (MB)},
        xmin=0,
        xmax=54,
        ymode=log,
        log basis y={10},
        legend columns=4,
        legend style={font=\footnotesize, at={(0.5,1.03)}, anchor=south, draw=none},
        grid=major
    ]
        \addplot+[mark=*] coordinates {(5.35,214.3) (6.15,411.9) (6.63,220.8) (7.25,397.0) (8.47,352.9) (9.97,350.7) (10.09,484.7) (10.19,348.2) (21.61,945.1) (33.55,975.4) (41.46,1366.4) (51.22,2287.0)};
        \addlegendentry{1 byte}
        \addplot+[mark=square*] coordinates {(5.35,262.3) (6.15,476.8) (6.63,221.9) (7.25,468.1) (8.47,365.7) (9.97,340.6) (10.09,597.6) (10.19,342.6) (21.61,1077.3) (33.55,961.6) (41.46,1374.7) (51.22,2533.6)};
        \addlegendentry{2 bytes}
        \addplot+[mark=triangle*] coordinates {(5.35,158.1) (6.15,221.3) (6.63,159.3) (7.25,246.8) (8.47,236.4) (9.97,252.5) (10.09,293.9) (10.19,256.8) (21.61,580.8) (33.55,727.3) (41.46,1010.4) (51.22,1518.4)};
        \addlegendentry{Bloc 1 byte}
        \addplot+[mark=diamond*] coordinates {(5.35,168.6) (6.15,393.3) (6.63,207.3) (7.25,458.2) (8.47,366.8) (9.97,331.9) (10.09,586.2) (10.19,337.9) (21.61,950.4) (33.55,935.5) (41.46,1335.6) (51.22,2449.5)};
        \addlegendentry{Bloc 2 bytes}
        \node[font=\tiny, rotate=45, anchor=west] at (axis cs:5.35,168.6) {xml};
        \node[font=\tiny, rotate=45, anchor=west] at (axis cs:6.15,393.3) {ooffice};
        \node[font=\tiny, rotate=45, anchor=west] at (axis cs:6.63,207.3) {reymont};
        \node[font=\tiny, rotate=45, anchor=west] at (axis cs:7.25,458.2) {sao};
        \node[font=\tiny, rotate=45, anchor=west] at (axis cs:8.47,366.8) {x-ray};
        \node[font=\tiny, rotate=45, anchor=west] at (axis cs:9.97,331.9) {mr};
        \node[font=\tiny, rotate=45, anchor=west] at (axis cs:10.09,586.2) {osdb};
        \node[font=\tiny, rotate=45, anchor=west] at (axis cs:10.19,337.9) {dickens};
        \node[font=\tiny, rotate=45, anchor=west] at (axis cs:21.61,950.4) {samba};
        \node[font=\tiny, rotate=45, anchor=west] at (axis cs:33.55,935.5) {nci};
        \node[font=\tiny, rotate=45, anchor=west] at (axis cs:41.46,1335.6) {webster};
        \node[font=\tiny, rotate=45, anchor=west] at (axis cs:51.22,2449.5) {mozilla};
    \end{axis}
    \end{tikzpicture}
\end{figure*}

La Taula~\ref{tab:modeavg} resumeix el compromís entre reducció i temps. Els
modes de 2 bytes comprimeixen millor de mitjana perquè codifiquen parelles de
bytes i poden capturar patrons que el mode d'1 byte no veu. El mode per blocs
de 2 bytes és lleugerament millor que el global de 2 bytes, però la diferència
és petita; per tant, el guany principal ve de la mida del símbol, no només de
dividir el fitxer en blocs. En canvi, el mode per blocs d'1 byte és el més
ràpid: construeix arbres més petits i, si un bloc no millora la mida final,
el pot guardar sense Huffman en lloc de forçar una codificació poc útil.

La Fig.~\ref{fig:mode_time_by_file} confirma que el temps creix sobretot amb
la mida de l'entrada. També mostra que els modes per blocs solen quedar per
sota dels globals en temps de compressió, especialment en fitxers grans com
\texttt{samba}, \texttt{webster} i \texttt{mozilla}.

%% ================================================================
\section{Conclusió i Opinió Personal}

En aquesta pràctica s'ha implementat una aplicació completa de compressió
Huffman amb interfície gràfica, línia de comandes, format d'arxiu propi,
diversos modes de compressió i una eina de benchmark. L'arquitectura separa
la vista, el controlador i el model, i dins el model reparteix les
responsabilitats entre anàlisi, selecció de mode, escriptura, lectura,
progrés i preparació de la informació visual.

Els resultats experimentals sobre Silesia mostren que el mode de 2 bytes i
el mode per blocs de 2 bytes són els que obtenen les reduccions més altes en
la majoria dels fitxers. També mostren que la cua de prioritat no afecta la
taxa de compressió quan el desempat és determinista, però sí el temps de
construcció de l'arbre. En aquest conjunt de dades, el heap binari és la
opció més eficient de les tres implementades.

Des d'un punt de vista personal, una de les parts més interessants ha estat
la visualització de l'arbre de Huffman. A més de ser una part visualment
atractiva de l'aplicació, ajuda a entendre d'un cop d'ull com es reparteixen
les freqüències i per què alguns símbols acaben amb codis més curts que
altres. També ha estat útil comparar modes sobre fitxers diferents, perquè
els resultats mostren que no hi ha una única opció òptima: el millor mode
depèn clarament del tipus de dades i de la seva distribució interna.

Les proves unitàries han estat especialment importants perquè el criteri de
correcció d'un compressor és estricte: després de comprimir i descomprimir,
el fitxer restaurat ha de coincidir byte a byte amb l'original. Per això les
proves combinen casos de compressió i descompressió completa, errors de format, arxius truncats,
determinisme de l'arbre i comprovacions del benchmark.

Com a treball futur, seria interessant ampliar el benchmark amb altres
tipus de fitxers i estudiar si es possible afegir altres variants de Huffman que
proporcionin una millor compressió per a certs arxius.

%% ================================================================
\begin{thebibliography}{9}

\bibitem{ref_gfg_huffman}
GeeksforGeeks, ``Huffman Coding | Greedy Algo-3,'' 2025. [Online]. Available:
\url{https://www.geeksforgeeks.org/dsa/huffman-coding-greedy-algo-3/}

\bibitem{ref_silesia}
S.~Deorowicz, ``Silesia compression data set,'' Silesian University of
Technology. [Online]. Available:
\url{https://sun.aei.polsl.pl/~sdeor/index.php?page=silesia}

\end{thebibliography}

\end{document}
